% Generated draft: requires mathematical and source review.
\documentclass[11pt]{article}
\usepackage[margin=25mm]{geometry}
\usepackage{fontspec}
\setmainfont{Latin Modern Roman}
\usepackage{amsmath,amssymb,mathtools,unicode-math}
\setmathfont{Latin Modern Math}
\usepackage{enumitem,needspace,xurl,hyperref}
\setlength{\parindent}{0pt}
\setlength{\parskip}{5pt}
\setlength{\emergencystretch}{3em}
\newcommand{\N}{\mathbb{N}}
\newcommand{\Z}{\mathbb{Z}}
\newcommand{\Q}{\mathbb{Q}}
\newcommand{\R}{\mathbb{R}}
\newcommand{\C}{\mathbb{C}}
\newcommand{\F}{\mathbb{F}}
\newcommand{\A}{\mathbb{A}}
\newcommand{\PP}{\mathbb{P}}
\newcommand{\E}{\mathbb{E}}
\newcommand{\eps}{\varepsilon}
\newcommand{\dd}{\,\mathrm{d}}
\begin{document}
\section*{UnsolvedMath ID: 1199}
\textbf{Source problem code:} catalogue entry
\subsection{Definitions and mathematical statement}
Let $G$ be a finite group and $p$ a prime. A $p$-subgroup is a subgroup whose order is a power of $p$, and its normalizer is
\[
N_G(Q)=\{g\in G:gQg^{-1}=Q\}.
\]
Write $\operatorname{Irr}(G)$ for the ordinary irreducible complex characters of $G$. A $p$-block $B$ is a block in modular representation theory at $p$; equivalently, blocks are associated with the primitive central idempotents of the group algebra over an algebraically closed field of characteristic $p$, and partition the ordinary irreducible characters via the usual modular block correspondence. Write $\operatorname{Irr}(B)$ for the ordinary irreducible characters belonging to $B$. For a positive integer $m$, let $v_p(m)$ be the exponent of $p$ in $m$. The $p$-defect of $\chi\in\operatorname{Irr}(G)$ is
\[
d_p(\chi)=v_p(|G|)-v_p(\chi(1)),
\]
and the character count at defect $d\geq0$ is
\[
k_d(B)=|\{\chi\in\operatorname{Irr}(B):d_p(\chi)=d\}|.
\]
Analogous counts can be formed for blocks of local normalizers $N_G(Q)$, using the relevant local group order to define character defect.

The supplied record attributes to Everett Dade (1992) a conjectural local-to-global relationship between these fixed-defect character counts. Its literal question asks only whether there is a relationship between block character counts of $G$ and local subgroups. The supplied literature triage explicitly says that the actual conjectures use precise alternating-sum formulas and come in several variants, including Dade's Projective Conjecture. It supplies neither the alternating-sum formula nor the admissible subgroup chains, local-block selection, signs, defect hypotheses, projective data, or character-triple conditions. Therefore no single Dade formula can be reconstructed from the displayed question without importing an unspecified variant. The mathematical target is to recover and state the chosen variant's precise alternating-sum identity, then establish it in its full stated domain; no identity or omitted hypothesis is invented here.

The registry labels the record open. The supplied triage dated 2026-08-17 labels it partially solved and reports the following scope of progress. Sp\"ath's reduction shows that the Character Triple Conjecture for quasisimple groups implies Dade's Projective Conjecture for all finite groups, with some quasisimple classes verified. Rossi proves the Character Triple Conjecture for $p$-solvable groups and gives a reduction of the general Character Triple Conjecture to quasisimple groups. Here a quasisimple group is a perfect group whose quotient by its centre is a nonabelian simple group; a $p$-solvable group has a normal series with factors of $p$-power order or order prime to $p$. The source reports no verified proof for all finite groups, and identifies verification of all relevant quasisimple covering-group conditions, together with selection of the exact Dade variant, as remaining tasks. These reduction statements do not identify all versions of Dade's conjecture with one another.
\subsection{Short English statement}
Dade's conjectures predict precise local formulas for counting ordinary irreducible characters of a fixed defect in a block of a finite group. The supplied question omits the formula and the choice among variants. Reported reduction theorems and special cases give progress but do not supply a general proof.
\subsection{Sources}
\begin{itemize}
\item Britta Sp\"ath, ``A reduction theorem for Dade's projective conjecture,'' Journal of the European Mathematical Society 19 (2017), 1071--1126: \url{https://ems.press/journals/jems/articles/14667}.
\item Damiano Rossi, ``Character Triple Conjecture for p-solvable groups,'' Journal of Algebra 595 (2022), 165--193: \url{https://doi.org/10.1016/j.jalgebra.2021.12.018}.
\item Damiano Rossi, ``A reduction theorem for the Character Triple Conjecture,'' arXiv:2402.10632 (2024): \url{https://arxiv.org/abs/2402.10632}.
\item Local registry data/unsolvedmath/problems.json, permanent catalogue entry, source code catalogue entry, ``Dade's Conjecture''; statement, background attribution to Everett Dade (1992), and literature triage dated 2026-08-17. The record gives no original 1992 bibliographic reference. Dataset attribution: Ulam AI and the UnsolvedMath Contributors.
\end{itemize}
\end{document}
